Extracting single-unit activity from raw, long-duration extracellular
recordings presents major challenges. Continuous neural recordings over weeks
or months, as in NaLoDuCo experiments, exhibit complex non-stationarities in
signal properties and produce massive datasets—often hundreds of terabytes in
size. These characteristics break many assumptions underlying conventional
spike sorting approaches and make manual curation infeasible.

To address these challenges our approach integrates experimental and
computational strategies. Experimentally, we will minimise electrode-tissue
movement by cementing probes to the skull, as in~\cite{schoonoverEtAl21}.
Computationally, we will create spike sorting strategies addressing
non-stationarities in NaLoDuCo recordings, distribute computation across
multiple nodes and use GPU acceleration to manage scale, and create automatic
curation methods tailored to non-stationarities in NaLoDuCo experimentation.

\subsubsubsection{Spike sorting approach for long-duration NaLoDuCo recordings}

In order to perform spike sorting on long-duration datasets, we will use an
overlapping windowing approach. Spike sorting will be done in short-duration
windows of a few hours (e.g., five hours), with a partial overlap between
consecutive windows (e.g., one hour).  We will first run spike sorting (e.g.,
Kilosort4) independently for each window. Next, we will use the overlapping
segment to compare the spike trains extracted from the two consecutive windows,
using existing methods implemented in SpikeInterface~\cite{buccinoEtAl20}. Assuming that
the spike sorter will find the same units and (most of) their spikes in that
period, unit identities will be tracked by finding the unit pairs on
consecutive sessions with the highest spike-train agreement in the overlapping
interval (we refer to this as consecutive tracking).

\subsubsubsection{Automatic quality control}

The long-duration of NaLoDuCo recordings and the large number of individual
spike sorting jobs makes manual inspection and curation unfeasible. We will
therefore integrate and evaluate the performance of two recently developed
automated tools to label spike sorted units: Bombcell\cite{fabreEtAl23} and
UnitRefine\cite{jainEtAl25}.  Both tools assign quality labels to individual
units (“good,” “multiunit activity” or “noise”)

After evaluating and fine-tuning the performance of Bombcell and UnitRefine on
individual spike sorting outputs, we will run automatic quality control
separately for each short-window spike sorting output. Using the new
visualisations we will develop, researchers will be able to explore quality
trends over timescales ranging from hours to weeks, detect periods of
instability, and automatically flag clusters exhibiting inconsistent or
transient quality metrics.

\subsubsubsection{Long-term tracking of units}

The extended duration of NaLoDuCo recordings introduces various forms of
non-stationarity, including tissue–electrode interactions, probe degradation,
neural plasticity, and brain state fluctuations. These factors can alter
neuronal waveforms over time and complicate spike sorting. To address this, we
will complement local unit tracking with long-term global tracking by extending
methods such as UnitMatch~\cite{vanBeestEtAl24} or the earth-moving-distance-based
approach~\cite{yuanEtAl24}. These extensions will maintain a catalogue of all
detected units, enabling reconciliation of unit identities across epochs where
waveforms shift or units transiently disappear.

\subsubsubsection{Spike sorting on the cloud/institutional HPC systems}

The very large size of NaLoDuCo recordings requires running spike sorting where
the data resides. We will build a web-based application to launch scalable
spike sorting jobs on multiple backends, including cloud providers and HPC
systems, and to visualise spike sorting results.

To run spike sorting at scale, we will extend an
\href{https://github.com/AllenNeuralDynamics/aind-ephys-pipeline}{existing
pipeline developed at AIND for spike sorting of Neuropixels data}. The pipeline
will be adapted for long-duration NaLoDuCo recordings to massively parallelize
all spike sorting jobs of short-duration windows, to automatically curate their
spike sorting results, to collect and compare the overlapping windows to track
unit identities, and to track long-term identities.

To visualise spike sorting results over extended periods will use the methods
described in Section Visual Exploration.

\subsubsubsection{Deliverables}

\begin{enumerate}

	\item Development of a spike sorting strategy for high-channel-count and long-duration recordings.

	\item Automated unit tracking and quality control across extended recordings.

	\item Web-based application to launch spike-sorting jobs on remote servers and to visualise long-duration spike sorting results.

	\item Spike sorting pipeline for cloud/HPC deployments for scalable and reproducible processing.

	\item Open-source release of all spike sorting tools, quality metrics, cloud deployment scripts, and documentation.

\end{enumerate}

